\spread{Can a writer learn the words around a fact?}
\noindent We held the visual reader fixed and trained two writers to turn its evidence into a complete answer. The same selected name could enter two different claims.
\platehere{Figure_173_writing_around_the_fact}{Supervising the surrounding text can recover a subject and lose part of a place name}{Rows show the actual question, selected evidence, human reference and three stopped replies. Panel a is the first gain over matched direct training in the English test order that fits the display. The earlier writer also gave the correct answer. Panel b is the corresponding first Chinese loss. The place is identified more precisely in the source and in the accepted human answer. Only spacing between original glyph cells changes for print.}
\begin{multicols}{2}
\section*{A controlled change to the training history}
Each writer receives a pictured question, the visual reader's selected strip and the full page. A recurrent network can draw a glyph from its image inventory, copy one that it sees or stop. Both fits start from the same earlier writer and keep the full-answer likelihood for every example. If an annotated value occurs once with surrounding words, the second fit also learns to write those words with the value replaced by one neutral cell in the preceding answer. The question and page still contain their evidence. The neutral cell is never scored as a required glyph. The control instead scores evenly spaced positions of the complete answer under its ordinary history.
\section*{What the learning signal preserves}
For reference $LFR$, with left context $L$, value $F$ and right context $R$, the auxiliary history before the $k$th right-hand glyph contains $L$, one gap and the preceding glyphs of $R$. The additional objective scores $|L|+|R|+1$ positions, including stop, in both arms. Lean checks this budget and the prefix identity. On a real 37-position English answer, independent single-step replay agreed with batched differentiation for all 40 parameter tensors in each arm (maximum absolute difference $7.2\times10^{-7}$).
\section*{Human and teacher supervision}
Each arm saw two passes over 133,048 complete reference wordings, with 26,365,654 ordinary and 4,660,116 auxiliary positions in 9,798 updates. The corpus includes human English statements, Chinese short answers, DuReader explanations, English Dolly responses and separately sourced bilingual Alpaca instruction responses. The two runs took 203 and 225 seconds on one RTX 4090 D, consuming 27.9 and 31.8 kJ of device energy. Their shared visual reader and input preparation have additional costs.
\end{multicols}
\newpage
\spread{Giving a writer more of the page}
\noindent Two questions were asked on each of 96 new passages per language. The test also included 128 human instruction responses absent from the pictured question and excerpt.
\platehere{Figure_174_complete_answers_during_training}{The fixed reader delivers the same facts to three writers}{Panel a counts stopped answer images exactly matching an accepted reference on the new passages. The earlier system passes only the chosen fact to its writer. Both new systems also pass the page. Panel b counts free replies on a separate, previously observed development cohort. Update zero uses the starting weights with the new inputs and inventory. Intermediate points use actual training snapshots.}
\begin{multicols}{2}
\section*{Where free answers deteriorated}
On English statements the earlier, direct and context writers gave 12, 6 and 5 exact answers out of 192. On Chinese short answers they gave 50, 40 and 38. The earlier model answered both questions correctly on two English and seven Chinese pages. The context fit did so on zero and five. Under a change to the evidence, both English answers were correct on three of 37 compatible edits for the earlier writer and one for each new writer. Even when supplied the human fact, the English counts dropped from 44 to 11 and 9. Improvements in individual printed sentences coexist with a loss in aggregate performance.
\section*{Source overlap and longer writing}
The fitted readers and writers shared no complete question, page or writer reference images with the reserved examples. A search for shared 64-glyph stretches identified four English statement questions and two Dolly examples with partially overlapping source text. Supplementary scores remove the affected families. All three models gave zero complete reference matches on 64 English and 64 Chinese instruction responses. Five English references exceeded the 1,024-decision budget. The context writer stopped on 37 English and 41 Chinese responses. A blinded local model accepted three of eight Chinese context replies and zero of eight English replies. It accepted every human answer and rejected every displaced answer in this small prose sample.
\section*{Recognizing the answer and writing it}
With two complete human answers supplied, we held the fact constant across the two questions and compared the two assignments by mean log likelihood per glyph. The context writer chose the human assignment on 95 of 96 English pages. Reversing the order of question glyphs reduced this to 50; duplicating one question produced 96 ties. Untrained Dice overlap of individual glyphs also matched 95 pairs; the two methods failed on different pages. Counting shared three-glyph runs matched 93. Thus aggregate accuracy on this test is attainable without learning sentence structure. On these same pages it freely answered both questions correctly zero times. Removing the rest of the page from this writer reduced exact replies from five to two in English and from 38 to 25 in Chinese. The earlier writer, trained on the fact strip alone, fell from 12 to one and from 50 to zero when given the whole page.
\section*{Can a wider search recover the missing answer?}
On the preceding 384 development questions we retained four live prefixes and scored every visited stop in the earlier 256-decision writer. Each additional glyph has conditional probability at most one. Hence a stopped reply outranking every discarded prefix is a global mode within the budget. At numerical precision this bound held for 167 English and all 192 Chinese questions. Of these modes, 161 and 149 respectively missed every accepted bitmap. When the read fact exactly matched the annotation, nine of 13 English modes still missed all references, whereas all 39 Chinese modes matched. Wider search recovered one English reference. A paraphrase can also miss an exact reference.
\end{multicols}
\newpage
\spread{When two questions share the same page}
\noindent We reserved further human-written question pairs and measured how often shared ink alone exposed the assignment of their answers.
\platehere{Figure_175_two_questions_about_one_company}{The same reader identifies a person in one answer and gives a person's name for a company in the other}{Both questions refer to the same pictured article about ASE Group. The first asks which company the group acquired in 2010. The human answer names Huanlong Dianqi. The model gives a person's name. The second asks who inspired the group's name. The model returns the accepted name of Master Hsing Yun. These are the first paired Chinese questions in the reserved overlap challenge with one correct free answer and outputs within the display limit. All image strips preserve the original glyph pixels.}
\begin{multicols}{2}
\section*{A surface control for the English result}
We scored each question and answer by their multiset overlap of single glyphs and of three-glyph runs. For each page we compared the sum of the two annotated similarities with the sum after exchanging the answers. With the preceding fits and test sets excluded, the remaining English source families yielded 85 new pairs. None met the condition that neither overlap measure favors the human pairing. On these pages both untrained measures chose the human assignment 84 times, and the retained writer chose it 83 times. The writer produced six of the 170 answers exactly and answered both questions on zero pages. High recognition accuracy here gives little leverage on whether a model has learned the roles expressed by the question.
\section*{Chinese questions without the overlap cue}
We reserved 48 Chinese pages on which neither measure favors the human assignment, then 96 pages from the remaining families. The retained writer assigned the references correctly on 20 of the 48 challenge pages, giving the same source evidence to both questions for scoring. Reversing the pictured questions gave 22. In free generation it matched 19 of 96 answers and answered both on zero pages. Supplying the human fact raised this to 94 of 96. On the 96 remaining pages it matched 36 of 192 on its own and 187 with the human fact. The overlap scores determine admission to the challenge. These rates describe the selected contrasts.
\end{multicols}
\newpage
